\documentclass[11pt,a4paper]{article}
\usepackage[margin=25mm]{geometry}
\usepackage[T1]{fontenc}
\usepackage{lmodern}
\usepackage{amsmath,amssymb,mathtools,bm}
\usepackage{microtype}
\usepackage{xcolor}
\usepackage{listings}
\usepackage[colorlinks=true,linkcolor=blue!55!black,citecolor=blue!55!black,urlcolor=blue!55!black]{hyperref}
\lstset{basicstyle=\ttfamily\small,breaklines=true,columns=fullflexible,keepspaces=true,frame=single}
\newcommand{\C}{\mathbb C}
\newcommand{\R}{\mathbb R}
\newcommand{\E}{\mathbb E}
\newcommand{\tr}{\operatorname{tr}}
\newcommand{\diag}{\operatorname{diag}}
\newcommand{\rank}{\operatorname{rank}}
\newcommand{\RePart}{\operatorname{Re}}
\newcommand{\norm}[1]{\left\lVert#1\right\rVert}
\title{Fisher Information for Pure-State Phase Retrieval\\\large Comparing phase screens through local reconstruction sensitivity}
\author{Study note for \texttt{random\_phase\_measurements}}
\date{October 6, 2026}
\begin{document}
\maketitle

\begin{abstract}
A normalized complex state in dimension $d$, modulo global phase, has
$2d-2$ real parameters. This note constructs local coordinates for those
parameters and derives the classical Fisher information for intensity
measurements, with the aim of comparing reconstruction quality as the
phase-screen parameters, especially the relative correlation distance $\zeta=r_c/w$, change.
The main analysis uses a fixed estimated signal power and background,
appropriate to the intended intense-laser experiment with negligible camera
cropping. It relates the inverse Fisher trace to local infidelity and
distinguishes measurement sensitivity from optimizer convergence. A detailed
unknown-power treatment is retained in Appendix~\ref{app:unknown-power}.
\end{abstract}

\section{The estimation problem}
Let $\{u_j\}_{j=1}^d$ be an orthonormal mode basis, and write the input field as
\begin{equation}
 U(\mathbf r)=\sum_{j=1}^d x_j u_j(\mathbf r),\qquad
 x=\sqrt{s}\,\psi,\qquad \psi^\dagger\psi=1,\quad s>0.
\end{equation}
The desired state is the ray represented by $\psi$: multiplying it by a
constant phase does not change the state. The brightness $s$ is also absent
from the normalized-state fidelity, but it can affect the measured counts.
It must therefore be distinguished from global phase.

For one phase screen, let $A\in\C^{M\times d}$ contain the propagated basis
fields sampled at the $M$ camera pixels. In the point-sampling model,
\begin{equation}
 q_i(\psi)=|(A\psi)_i|^2,\qquad
 y_i\overset{\mathrm{ind}}\sim\operatorname{Poisson}(\lambda_i),\qquad
 \lambda_i=sq_i(\psi)+b_i, \label{eq:model}
\end{equation}
where $b_i$ is the known mean background in the same count units.
The Poisson phase-retrieval objective of Li, Lange, and Fessler is
$\sum_i[\lambda_i-y_i\log\lambda_i]$, up to constants independent of the
parameters \cite{li2022}. Here $s$ is a scale parameter; it equals the
expected detected signal count only when $\sum_iq_i=1$.

For finite pixel integration, replace $q_i$ by
$\psi^\dagger E_i\psi$, where $E_i$ integrates the propagated mode products
over pixel $i$. The rank-one expression in \eqref{eq:model} is its
point-sampling approximation. The camera-POVM viewpoint and finite-subspace
tomography are developed in Ref.~\cite{oliveira2025}.

For all Hermite--Gaussian modes of total order at most $L$,
\begin{equation}
 d=\frac{(L+1)(L+2)}{2}.
\end{equation}
At $L=4$, $d=15$: pure-state estimation needs 28 real state parameters,
whereas unrestricted density-matrix tomography needs 224.

\subsection{Experimental regime and power convention}
The objective is to compare phase screens, mainly as their relative correlation distance
$\zeta=r_c/w$ changes, rather than to study uncertainty in optical power.
The intended experiment uses intense laser illumination, an independently
estimated background, and camera coverage with no substantial cropping.
We therefore treat the estimated signal power and background as fixed in
the main Fisher calculation. Their estimation uncertainty is neglected in
this approximation; the background is not assumed to be physically zero.

If the calibrated propagation approximately preserves power on the mode
subspace, $A_\zeta^\dagger A_\zeta\approx I$, a per-image estimate is
\begin{equation}
 \widehat s=\sum_i(y_i-\widehat b_i),\qquad
 \widehat{\sqrt{s}}=\sqrt{\widehat s}.
 \label{eq:estimated-power}
\end{equation}
Here $\widehat s$ is required to be positive. The sum uses the signed
background-subtracted values, without clipping individual pixels, since
clipping would bias the total upwards. If the common power convention is
$A_\zeta^\dagger A_\zeta\approx\alpha I$, divide the sum by $\alpha$.
Column normalization alone does not establish either identity; the basis
must also remain approximately orthogonal on the detector grid.

Background subtraction in \eqref{eq:estimated-power} is used only to
estimate total signal power. The Poisson likelihood still receives the
original images and uses $\lambda_i=sq_i+\widehat b_i$. It does not receive
a clipped, background-subtracted image. For the remainder of the main
analysis, $s$ and $b$ denote fixed values, including plug-in estimates.
The existing unconstrained solver may continue to fit the norm of $x$
implicitly and normalize its output; this diagnostic convention does not
require changing the solver.

For controlled comparisons between $\zeta$ values, use a common reference
signal power $s_{\mathrm{ref}}$, background model, and set of test states.
This isolates the change in the measurement from incidental brightness
variation. Per-image power estimates remain useful for interpreting actual
experimental runs, but their Fisher values correspond to those runs' own
resources. If backgrounds vary with exposure, include that variation when
assessing the actual acquisition conditions and state it separately from
the comparison at common power.

In the background-free, power-preserving limit, state and power are exactly
Fisher-orthogonal, so treating power as fixed causes no loss of state
information relative to fitting it. Weak background motivates using this
as an approximation in the intense regime. The exact correction for unknown
power, including its assumptions and numerical implementation, is developed
in Appendix~\ref{app:unknown-power}.

\section{Characterizing phase-screen scale and strength}\label{sec:screen-parameters}
The phase-screen parameters should describe two separate physical properties:
the spatial scale of the phase variations and their statistical strength.
For the main comparison, the strength is held fixed while the spatial
scale is varied. This section describes the revised generator defaults:
a $1/e$ correlation distance $r_c$ and spectral amplitude $2\pi^2$.

\subsection{Spatial correlation distance and its configuration convention}
The generator in \texttt{acquisition/generate\_phase\_masks.py} uses
\begin{equation}
 S(\mathbf q)=\frac{\mathcal A r_c^2}{4\pi}
 \exp\!\left[-\frac{r_c^2(q_x^2+q_y^2)}4\right],
 \qquad \mathcal A=\texttt{amplitude}=2\pi^2\quad\text{by default}.
 \label{eq:screen-spectrum}
\end{equation}
Its frequency grid uses angular spatial frequency:
$q_x=2\pi k_x/(N_x\,dx)$ and $q_y=2\pi k_y/(N_y\,dy)$, with the usual
FFT ordering of signed indices. For $dx=dy=1$, $r_c$ is in pixels.
There is no extra factor of $2\pi$ in the conversion below.

Let $g(\mathbf r)$ denote the zero-mean Gaussian phase field before wrapping,
with ensemble RMS $\phi_{\mathrm{rms}}=\sqrt{\E[g^2]}$.
By the covariance--spectrum Fourier relation, the continuous, large-grid
approximation gives
\begin{equation}
 C_g(r)=\E[g(\mathbf r)g(\mathbf0)]
 =\phi_{\mathrm{rms}}^2\exp\!\left(-\frac{r^2}{r_c^2}\right),
 \qquad \boxed{C_g(r_c)/C_g(0)=e^{-1}},\quad \phi_{\mathrm{rms}}^2=C_g(0).
 \label{eq:phase-covariance}
\end{equation}
Thus $r_c$ is the $1/e$ correlation distance of the unwrapped phase.
In the standard squared-exponential covariance convention
$\phi_{\mathrm{rms}}^2\exp[-r^2/(2\ell^2)]$ \cite{rasmussen2006}, $\ell=r_c/\sqrt2$.
The earlier spectral-width convention is related by
\begin{equation}
 \sigma=\frac{\sqrt2}{r_c},\qquad \ell=\frac1\sigma.
\end{equation}
This is a reparameterization of the same Gaussian spectral family.

For a beam waist $w$, define the dimensionless parameter
\begin{equation}
 \boxed{\zeta=\frac{r_c}{w}.}
\end{equation}
Correlation distances are specified in units of the beam waist, as the
ratios $\zeta=r_c/w$. The script converts them to grid units by multiplying
by the waist before generating the screens and saving the lengths in the
HDF5 \texttt{correlation\_lengths} dataset. Dividing these saved lengths by
the waist recovers the specified ratios. Larger $\zeta$ produces
smoother screens; smaller $\zeta$ produces shorter variations and larger
phase gradients. Plots may use $\zeta$ directly, or $1/\zeta=w/r_c$ if an
axis increasing with roughness is preferred. Throughout the measurement
comparisons below, subscripts $\zeta$ label the screen scale; individual
realizations at the same scale still have different measurement matrices.

\subsection{What the amplitude parameter actually controls}
The normalization in \eqref{eq:screen-spectrum} gives
$\int_{\R^2}S(\mathbf q)\,d^2q=\mathcal A$. In the current implementation,
independent unit-variance complex Gaussian coefficients are multiplied by
$\sqrt{S\,\Delta q_x\Delta q_y}$, Fourier transformed with no inverse
normalization, and then their real part is taken. Each complex Gaussian
coefficient has real and imaginary variances $1/2$. It follows that
\begin{equation}
 \phi_{\mathrm{rms,grid}}^2=\frac12\sum_{k_x,k_y}
 S(q_x,q_y)\,\Delta q_x\Delta q_y
 \approx\frac{\mathcal A}{2},\qquad
 \boxed{\phi_{\mathrm{rms}}=\sqrt{\E[g^2]}
 \approx\sqrt{\frac{\mathcal A}{2}}.}
 \label{eq:phase-rms}
\end{equation}
Here $\phi_{\mathrm{rms}}$ denotes the underlying zero-mean, unwrapped
ensemble RMS phase, in radians. It is not the RMS of the displayed wrapped
angles. The factor $1/2$ belongs to this particular complex-noise/real-part
construction and should be revisited if the generator is changed.

The generator returns $\phi=\operatorname{Arg}(e^{ig})$, implemented as
\texttt{angle(exp(1j*g))}. Its stored angles lie in the principal range
$[-\pi,\pi]$, regardless of $\mathcal A$, and $e^{i\phi}=e^{ig}$ exactly.
This range is a representation convention, not a measure of phase strength.

The spectral amplitude and corresponding strengths are:
\begin{center}
\begin{tabular}{c|c}
Spectral parameter $\mathcal A$ & Approximate unwrapped $\phi_{\mathrm{rms}}$\\\hline
$\pi$ (earlier setting) & $\sqrt{\pi/2}\simeq1.25$ rad\\
$\pi^2$ & $\pi/\sqrt2\simeq2.22$ rad\\
$2\pi^2$ (implemented default) & $\pi\simeq3.14$ rad\\
$\pi^2/2$ (optional weaker comparison) & $\pi/2\simeq1.57$ rad
\end{tabular}
\end{center}
Fixing $\phi_{\mathrm{rms}}=\pi$ is a deliberate strong-modulation baseline,
not a claim of an optimal strength or of a hard unwrapped phase limit.
The main sweep varies $\zeta=r_c/w$ at this fixed strength. A smaller
strength, such as $\pi/2$, can be tested on a subset to check whether the
qualitative conclusions depend on the strength choice. The current acquisition
script uses the generator's fixed amplitude default; an amplitude entry in
the configuration is not used by this script.

\subsection{Correlation of the phase factor actually applied to light}
The optical transformation depends on $P=e^{i\phi}$, rather than on the
choice of representative angle. Gaussian phase differences satisfy
\begin{equation}
 D_g(r)=\E[(g(\mathbf r)-g(\mathbf0))^2]
 =2\phi_{\mathrm{rms}}^2\left(1-e^{-r^2/r_c^2}\right).
\end{equation}
For a zero-mean Gaussian variable $Z$, $\E[e^{iZ}]=e^{-\operatorname{Var}(Z)/2}$.
Applying this identity to the phase difference gives
\begin{equation}
 R_P(r)=\E[P(\mathbf r)P(\mathbf0)^*]
 =\exp\!\left[-\phi_{\mathrm{rms}}^2\left(1-e^{-r^2/r_c^2}\right)\right].
 \label{eq:phase-factor-correlation}
\end{equation}
This is the phase-factor correlation, normalized by $R_P(0)=1$, without
subtracting its mean. It approaches $e^{-\phi_{\mathrm{rms}}^2}$ at large separation.
For $g$, $|\E[e^{ig}]|=e^{-\phi_{\mathrm{rms}}^2/2}$, so that limit is the squared coherent mean.
For the default $\phi_{\mathrm{rms}}\approx\pi$, the mean magnitude is about $0.0072$ and the
wrapped one-point phase distribution is nearly uniform. This does not
mean that nearby pixels are independent.

For $r\ll r_c$,
\begin{equation}
 R_P(r)\approx\exp\!\left(-\frac{\phi_{\mathrm{rms}}^2 r^2}{r_c^2}\right).
\end{equation}
Consequently, for strong modulation its short-distance length scale is
approximately $r_c/\phi_{\mathrm{rms}}$, when defined at correlation $1/e$. The underlying phase correlation
length and the optical phase-factor correlation length are distinct;
changing strength changes the latter even when $r_c$ is fixed. At fixed
$\phi_{\mathrm{rms}}$, the mean-square phase gradient is
\begin{equation}
 \E[|\nabla g|^2]=\frac{4\phi_{\mathrm{rms}}^2}{r_c^2}=2\phi_{\mathrm{rms}}^2\sigma^2,
\end{equation}
which makes the increase in gradient roughness as $r_c$ decreases explicit.
Gaussian screens remain mathematically smooth; ``rougher'' here means
shorter spatial variations and larger typical gradients, not a change in
differentiability class.

\subsection{Finite-grid and empirical characterization}
The Fourier construction is periodic, and its exact covariance is a
finite Fourier sum rather than the infinite-plane Gaussian in
\eqref{eq:phase-covariance}. The approximation requires resolving the
spectral peak and avoiding significant truncation at the Nyquist frequency.
The grid formula in \eqref{eq:phase-rms} provides the actual ensemble
variance for the chosen discretization.

Individual realizations, and especially their illuminated portions, need
not have exactly the prescribed RMS. The zero-frequency component also
contributes an irrelevant random global phase; removing the spatial mean
removes its contribution from a screen's realized variance. For reported
strength, distinguish the prescribed ensemble RMS from the measured RMS
after removing that mean. Use multiple realizations rather than rescaling
every screen by its extrema, which would change the statistical ensemble.

To check spatial scale empirically, retain the unwrapped field and examine
its covariance or structure function. With only wrapped angles available,
use $P(\mathbf r)P(\mathbf r+\Delta\mathbf r)^*$ and average over positions
and realizations, comparing with \eqref{eq:phase-factor-correlation}.
Ordinary differences of wrapped angles can introduce artificial jumps at
the branch cut. Expressing scale as $r_c/w$ and strength as
$\phi_{\mathrm{rms}}$ makes the design reproducible and separates the two
controls relevant to the rank and Fisher comparisons.

\section{Removing normalization and global phase}
Using the real and imaginary parts of all $d$ coefficients gives $2d$
real coordinates. Normalization removes one degree of freedom; identifying
$\psi$ with $e^{i\phi}\psi$ removes another. The physical parameter space is
complex projective space $\mathbb{CP}^{d-1}$.

Fix a normalized reference state $\psi$, at which the information is to be
evaluated. Choose $Q\in\C^{d\times(d-1)}$ satisfying
\begin{equation}
 Q^\dagger Q=I_{d-1},\qquad Q^\dagger\psi=0.
 \label{eq:Q}
\end{equation}
An SVD nullspace of $\psi^\dagger$ supplies such a basis. For real
$u,v\in\R^{d-1}$, define the local chart
\begin{equation}
 \psi(\theta)=\frac{\psi+Q(u+iv)}{\sqrt{1+\norm{u}^2+\norm{v}^2}},
 \qquad \theta=\begin{pmatrix}u\\v\end{pmatrix}\in\R^{2d-2}.
 \label{eq:chart}
\end{equation}
It is normalized exactly and automatically fixes the global phase, because
its overlap with the reference state is real and positive:
\begin{equation}
 \psi^\dagger\psi(\theta)
 =\frac{1}{\sqrt{1+\norm{u}^2+\norm{v}^2}}>0.
\end{equation}
This follows directly from $\psi^\dagger Q=0$; no additional constraint
needs to be imposed. It describes a neighborhood of the reference
ray; no single global chart is needed for a local Fisher calculation.

At $\theta=0$, introduce
\begin{equation}
 H=[\,Q\quad iQ\,],\qquad \delta\psi=H\delta\theta.
\end{equation}
Then
\begin{equation}
 \psi^\dagger\delta\psi=0,\qquad
 \norm{\delta\psi}^2=\norm{\delta\theta}^2,\qquad
 \RePart(H^\dagger H)=I_{2d-2}.
 \label{eq:metric}
\end{equation}
Thus $H$ is an orthonormal tangent basis over the \emph{real} numbers.
The excluded complex line contains both the radial direction $\psi$ and
the phase direction $i\psi$. This construction is consistent with the
local pure-state coordinates used in Fisher-symmetric tomography
\cite{li2016}.

\section{Intensity Jacobian and Poisson Fisher information}
Set $w=A\psi$ and $B=AH=[\,AQ\quad iAQ\,]$. Differentiating
$q_i=w_i^*w_i$ gives
\begin{equation}
 J_{ik}=\left.\frac{\partial q_i}{\partial\theta_k}\right|_0
 =2\RePart(w_i^*B_{ik}),\qquad
 J\in\R^{M\times(2d-2)}. \label{eq:J}
\end{equation}
For general pixel operators, the corresponding formula is
$J_{ik}=2\RePart(\psi^\dagger E_i H_{:k})$.

For arbitrary real parameters $\eta$, the log-likelihood score is
\begin{equation}
 \partial_a\ell=\sum_i\left(\frac{y_i}{\lambda_i}-1\right)
 \partial_a\lambda_i.
\end{equation}
Using independence and $\operatorname{Var}(y_i)=\lambda_i$ yields the
\emph{expected classical Fisher information}
\begin{equation}
 F_{ab}=\E[(\partial_a\ell)(\partial_b\ell)]
 =\sum_i\frac{\partial_a\lambda_i\,\partial_b\lambda_i}{\lambda_i}.
 \label{eq:fisher}
\end{equation}
For known $s$ and $b$, let $W=\diag(1/\lambda_i)$. Then
\begin{equation}
 \boxed{F_{\mathrm{pure}}=F_{\theta\theta}=s^2J^{\mathsf T}WJ.} \label{eq:known}
\end{equation}
All matrices in this expression except the intermediate optical fields are
real. The displayed formulas assume strictly positive $\lambda_i$.

\subsection{Why the optimizer's Fisher expression is different}
The WF step-size rule in Ref.~\cite{li2022}, also used by the local
\texttt{PoissonPhaseRetrieval} package, employs the complex matrix
\begin{equation}
 A^\dagger\diag\!\left(\frac{4|(Ax)_i|^2}{|(Ax)_i|^2+b_i}\right)A.
 \label{eq:wf}
\end{equation}
This is the complex score-based quantity used in that update rule. It is
not the real-parameter Fisher matrix for normalized states modulo phase.
For example, a perturbation $\delta x=ix$ leaves every intensity unchanged:
its physical Fisher information is zero. The quadratic form of
\eqref{eq:wf} along $ix$ is generally nonzero. Therefore one should derive
the diagnostic from the real intensity Jacobian, rather than invert the
step-size matrix. This distinction does not invalidate the WF update rule.

\section{Global uniqueness and the proposed results}\label{sec:uniqueness}
The motivation for phase modulation is that a direct image can fail to
identify even a pure state. Local Fisher information and global uniqueness
must therefore be distinguished before interpreting a reconstruction result.

\subsection{The order-1 direct-image ambiguity}
Choose real order-1 basis fields $h_1=\mathrm{HG}_{10}$ and
$h_2=\mathrm{HG}_{01}$ in the observation plane. For
$U=a h_1+b h_2$,
\begin{equation}
 |U|^2=|a|^2h_1^2+|b|^2h_2^2
 +2\RePart(a^*b)h_1h_2.
\end{equation}
The image does not depend on the imaginary coherence. In Bloch coordinates
$\rho=(I+r_x\sigma_x+r_y\sigma_y+r_z\sigma_z)/2$, it is
\begin{equation}
 I(\mathbf r)=\frac{1+r_z}{2}h_1^2+
 \frac{1-r_z}{2}h_2^2+r_xh_1h_2.
\end{equation}
The two LG states proportional to $h_1\pm i h_2$ are opposite poles of
the unobserved $r_y$ axis. Every line segment in the Bloch ball parallel
to that axis consists of mixed states with the same image. A generic such
line intersects the pure-state sphere twice, giving indistinguishable
pure states related by $r_y\mapsto-r_y$.

Nevertheless, the pure-state tangent Fisher matrix can have full rank at
$r_y\ne0$: within a fixed hemisphere, the measured coordinates $(r_x,r_z)$
locally specify a point on the sphere. The reflected state remains globally
indistinguishable. At $r_y=0$, the first-order tangent sensitivity also loses
a direction. Thus full tangent Fisher rank at a sampled state is not a
certificate of global pure-state uniqueness.

\subsection{Full-tomography rank as a sufficient uniqueness test}
For each phase screen, let $T_\zeta$ be the traceless part of the
full-tomography measurement matrix, so that the signal measurements have
an affine dependence on Bloch coordinates:
\begin{equation}
 q(\rho)=c_\zeta+T_\zeta t,\qquad
 q(\rho_1)-q(\rho_2)=T_\zeta(t_1-t_2).
\end{equation}
At fixed nonzero power and known background, equality of Poisson means
is equivalent to equality of these signal measurements. Therefore
\begin{equation}
 \rank(T_\zeta)=d^2-1
 \quad\Longrightarrow\quad
 \text{global uniqueness for all density matrices, hence all pure states}.
 \label{eq:full-rank-sufficient}
\end{equation}
If $T_\zeta$ is deficient, a nonzero kernel direction generates
indistinguishable physical density matrices around the maximally mixed
state. It does not necessarily generate indistinguishable pure states.
Full-tomography rank is therefore sufficient, but not necessary, for the
pure-state objective.

For strictly positive outcome probabilities, the full-tomography Fisher
matrix has the same rank as $T_\zeta$, since its diagonal weighting is
invertible. Its rank consequently has a global implication because the
measurement map is affine, even though the Fisher matrix itself describes
local uncertainty. Its eigenvalues and inverse trace remain state dependent.

\subsection{The exact pure-state criterion, when needed}
Define the linear measurement map on Hermitian matrices by
\begin{equation}
 \mathcal L_\zeta(X)_i=\tr(E_{\sigma,i}X).
\end{equation}
Two normalized pure states have the same image exactly when
$\mathcal L_\zeta(\psi\psi^\dagger-\phi\phi^\dagger)=0$.
Their difference is Hermitian, traceless, and has rank at most two.
Conversely, a nonzero traceless Hermitian matrix of rank at most two has
spectral form $a(uu^\dagger-vv^\dagger)$, where $a\ne0$ and $u,v$ are
orthonormal. Hence
\begin{equation}
\boxed{\ker\mathcal L_\zeta\ \cap\
 \{X=X^\dagger:\ \tr X=0,\ \rank X\le2\}=\{0\}.}
\end{equation}
This condition is equivalent to global uniqueness among normalized pure states.
This is the normalized-state specialization of low-rank-kernel criteria
in tomography with prior information and complex phase retrieval
\cite{heinosaari2013,conca2015}.

An optional global separation diagnostic is
\begin{equation}
 \beta_\zeta=\min_{\substack{X=X^\dagger,\ \tr X=0\\
 \rank X\le2,\ \norm{X}_F=1}}\norm{\mathcal L_\zeta(X)}_2.
\end{equation}
The constraint set is compact. Consequently, $\beta_\zeta>0$ is equivalent
to global uniqueness, and zero indicates an exact ambiguity. It can be
searched over $X=(uu^\dagger-vv^\dagger)/\sqrt2$ with orthonormal $u,v$.
A local numerical search that finds a collision provides evidence of
nonuniqueness, subject to residual tolerance; failing to find one does not
certify a positive global minimum. This optimization is not required for
the proposed main results wherever \eqref{eq:full-rank-sufficient} applies.

\subsection{Two complementary results versus screen roughness}
The proposed presentation consists of two separate results:
\begin{enumerate}
 \item \textbf{Full-tomography rank versus $\zeta$.}
 Plot $\rank(T_\zeta)$, or its normalized value
 $\rank(T_\zeta)/(d^2-1)$, and identify the sampled region where full
 column rank is obtained. For order up to 4, the target rank is 224.
 This establishes a sufficient global uniqueness criterion for the modeled
 measurement. Use the traceless matrix $T_\zeta$, not the complex optical
 matrix $A_\zeta$ and not the matrix including the identity column.
 In \texttt{QuantumMeasurements}, it is obtained with
 \texttt{get\_traceless\_part(measurement\_matrix)}.
 \item \textbf{Pure-state inverse Fisher trace versus $\zeta$.}
 Evaluate $\tr(F_{\mathrm{pure}}(\zeta;\psi)^{-1})$ at common signal power
 and background over the same normalized test states. Report a central
 statistic and spread across states and screen realizations. For order
 up to 4, this Fisher matrix has dimension 28. Lower inverse trace means
 a better local precision benchmark for normalized-state reconstruction.
 The metric remains local even in a region where global uniqueness has
 already been established by the first result.
\end{enumerate}

Numerical rank requires an explicit tolerance, for example counting
singular values larger than $\epsilon_{\mathrm{rel}}$ times the largest
singular value. State the grid, model, arithmetic precision, and tolerance,
and examine the smallest singular values near the apparent transition.
A weak phase screen can lift an exact degeneracy by an extremely small
amount: an observed onset of numerical full rank need not be a sharp
mathematical threshold in $\zeta$. Repeat the calculation for several
screen realizations; a region found on a sampled grid is not a guarantee
for every screen or every intermediate value of $\zeta$.

The second result shows precision within the identifiable region, which
rank alone cannot quantify. It can also be evaluated outside that region,
but full local Fisher rank there must not be presented as global uniqueness.
Neither roughness nor full rank guarantees monotonic improvement in the
inverse trace, and neither diagnostic guarantees convergence of a chosen
phase-retrieval algorithm. Experimental reconstruction fidelity remains
an additional test of the complete inference procedure.

\section{Connection to full tomography and fidelity}
Let $\{\omega_a\}_{a=1}^{d^2-1}$ be an orthonormal traceless Hermitian basis
and $t_a=\tr(\rho\omega_a)$. For a normalized POVM, write
$p=c+Tt$. The matrix returned by \texttt{QuantumMeasurements.fisher} for
an ordinary measurement matrix is
\begin{equation}
 F_{\mathrm{tom}}=T^{\mathsf T}\diag(1/p_i)T.
\end{equation}
It is per-detection information in all $d^2-1$ density-matrix coordinates.
Passing a pure state as its evaluation point does not change that dimension.

One can restrict this existing matrix to the pure-state tangent space.
At $\rho=\psi\psi^\dagger$, define
\begin{equation}
 R_{ak}=\frac{\partial t_a}{\partial\theta_k}
 =\tr\!\left[\omega_a(H_{:k}\psi^\dagger+\psi H_{:k}^\dagger)\right]
 =2\RePart(\psi^\dagger\omega_aH_{:k}).
\end{equation}
The chain rule gives
\begin{equation}
 F_{\mathrm{pure},1}=R^{\mathsf T}F_{\mathrm{tom}}R.
\end{equation}
For the same normalized measurement model, this equals the direct
probability-Jacobian calculation, with $p_i=q_i/\sum_jq_j$ and
$J_p=\partial p/\partial\theta$ for the conditional, background-free
model. This per-detection formula is distinct from the fixed-power Poisson
formula with additive background. Building the Jacobian directly avoids constructing
the much larger full-tomography matrix. For orthonormal $\omega_a$,
$R^{\mathsf T}R=2I$: this factor matters when comparing coordinate conventions.

Define fidelity by $\mathcal F(\psi,\phi)=|\psi^\dagger\phi|^2$.
The chart \eqref{eq:chart} gives the exact identity
\begin{equation}
 1-\mathcal F(\psi,\psi(\theta))
 =\frac{\norm{\theta}^2}{1+\norm{\theta}^2}
 =\norm{\theta}^2+O(\norm{\theta}^4).
\end{equation}
For a regular, locally unbiased estimator, the Cram\'er--Rao inequality
implies $\operatorname{Cov}(\widehat\theta)\succeq F_{\mathrm{pure}}^{-1}$.
Consequently,
\begin{equation}
 \E[1-\mathcal F]\ \gtrsim\ \tr(F_{\mathrm{pure}}^{-1})
 \label{eq:infidelity}
\end{equation}
is a \emph{local, small-error benchmark}, not an exact finite-sample
inequality for infidelity. For fixed $N$ and per-photon $F_1$, the benchmark
is $\tr(F_1^{-1})/N$. Asymptotically efficient estimators can attain the
leading term under suitable regularity conditions. Bias, model mismatch,
and convergence to another solution can invalidate that prediction.

The trace is meaningful here because the coordinates are orthonormal for
the local fidelity metric. In general coordinates $\xi$, if
$1-\mathcal F\simeq\delta\xi^{\mathsf T}M\delta\xi$, the corresponding
quantity is $\tr(MF_\xi^{-1})$. Under a change of $Q$ satisfying
\eqref{eq:Q}, the induced real coordinate change is orthogonal, so the
Fisher eigenvalues and inverse trace above are unchanged.

\section{What to compute and how to interpret it}\label{sec:diagnostics}
Define the weighted state Jacobian
\begin{equation}
 G=\diag(\lambda^{-1/2})sJ,\qquad F_{\mathrm{pure}}=G^{\mathsf T}G.
\end{equation}
This is a Gram matrix of the weighted state derivatives. Prefer computing
an SVD of $G$ directly: forming its Gram matrix squares the condition number
and can obscure small singular values. Let
$\tau_1\ge\cdots\ge\tau_{2d-2}$ be the singular values of $G$; the symbol
$\tau$ distinguishes them from the phase-screen parameter $\zeta$. Then
\begin{equation}
 \rank(F_{\mathrm{pure}})=\rank(G),\qquad
 \kappa_2(F_{\mathrm{pure}})=\left(\frac{\tau_1}{\tau_{2d-2}}\right)^2,
 \qquad \tr(F_{\mathrm{pure}}^{-1})=\sum_k\tau_k^{-2}.
\end{equation}
Use a stated numerical rank tolerance. A missing physical direction makes
the ordinary inverse-trace criterion divergent; using a pseudoinverse and
silently discarding that direction would make the measurement look better
than it is. A singular Fisher matrix means failure of regular first-order
identifiability; higher-order identifiability requires separate analysis.

Full rank is a local property. It does not exclude a distant state with the
same intensities, and it does not ensure that a particular initialization
converges to the desired solution. Condition number measures anisotropy,
not absolute precision: $\epsilon I$ is perfectly conditioned but carries
little information when $\epsilon$ is small. Report the inverse trace and
smallest eigenvalue as well, with a fixed resource convention (incident
brightness or detected photon number).

For each tested $\zeta$, screen realization, and normalized input state,
construct $A_\zeta$, the state tangent basis, and $J_\zeta$. At common
power and background, the comparison uses
\begin{equation}
 F_{\mathrm{pure}}(\zeta;\psi)
 =s_{\mathrm{ref}}^2J_\zeta^{\mathsf T}
 \diag\!\left(\frac{1}{s_{\mathrm{ref}}|(A_\zeta\psi)_i|^2+b_i}\right)
 J_\zeta.
\end{equation}
The diagonal entries are understood pixel by pixel. Evaluate the same
collection of Haar-sampled normalized states for all screens. Report the
fraction of deficient cases, median inverse trace, an upper quantile or
worst observed value, smallest eigenvalue, and condition number. Include
several screen realizations at each $\zeta$ to distinguish systematic
parameter effects from one favorable or unfavorable realization.

Compare these statistical diagnostics with experimental fidelities and
noiseless recovery from multiple initializations. Fisher information
predicts local sensitivity to measurement noise near the evaluated state;
it does not predict whether spectral initialization reaches the correct
optimization basin. A screen can have good local information and still
produce poor reconstruction with a particular optimizer. Thus fidelity
versus $\zeta$ reflects both the measurement design and solver behavior.

Zero-probability pixels need special care: the regular formulas can contain
$0/0$, and their limits may depend on the approach direction. A physical
positive background removes this particular singularity. Adding an arbitrary
numerical floor changes the model and should be disclosed. Likewise, raw
camera values require an appropriate gain/noise model; saturation and clipped
background subtraction are not independent Poisson counting processes.

\section{A compact Julia implementation}
The following computes the Poisson state information and its weighted
Jacobian for fixed estimated power and background. An optional argument
enables the unknown-power correction from Appendix~\ref{app:unknown-power}. It assumes
an orthonormal coefficient basis and strictly positive expected pixel means.
It accepts a scalar background or a vector of length $M$.

\begin{lstlisting}
using LinearAlgebra

function pure_state_fisher(A, psi; scale=1.0, background=0.0,
                          unknown_scale=false)
    scale > 0 || throw(ArgumentError("scale must be positive"))
    norm(psi) > 0 || throw(ArgumentError("state must be nonzero"))
    p = normalize(complex.(psi))
    Q = nullspace(reshape(conj.(p), 1, length(p)))
    H = hcat(Q, im * Q)
    w = A * p
    q = abs2.(w)
    J = 2 .* real.(conj.(w) .* (A * H))
    lambda = scale .* q .+ background
    all(isfinite, lambda) && all(>(0), lambda) ||
        throw(ArgumentError("expected means must be finite and positive"))
    G = (scale .* J) ./ sqrt.(lambda)
    if unknown_scale
        h = q ./ sqrt.(lambda)
        dot(h, h) > 0 || throw(ArgumentError("unobservable scale"))
        G -= h * ((h' * G) / dot(h, h))
    end
    return Symmetric(G' * G), G
end

# Example: use the true prepared coefficient vector to assess a screen.
# F, G = pure_state_fisher(A, coefficient;
#     scale=reference_signal_power, background=background_counts)
# singular_values = svdvals(G)
# Check numerical rank before computing:
# inverse_trace = sum(inv(abs2(t)) for t in singular_values)
\end{lstlisting}

The reference state is held fixed while constructing $Q$ and differentiating
the chart. For real data it may be replaced by an estimated state, producing
a plug-in local diagnostic. With the default \texttt{unknown\_scale=false}, the function returns
$F_{\mathrm{pure}}$ for the main fixed-power approximation. Setting it to
\texttt{true} instead returns $F_{\mathrm{eff}}$ from the appendix.
The function evaluates information only; it does not modify the
phase-retrieval optimizer.

\clearpage
\appendix
\section{Unknown power as a nuisance parameter}\label{app:unknown-power}
When $s$ is fitted rather than known, the joint Jacobian for $(\theta,s)$ is
\begin{equation}
 K=[\,sJ\quad q\,],\qquad F_{\mathrm{joint}}=K^{\mathsf T}WK
 =\begin{pmatrix}
 F_{\theta\theta}&F_{\theta s}\\
 F_{s\theta}&F_{ss}
 \end{pmatrix},
\end{equation}
with $F_{\theta s}=sJ^{\mathsf T}Wq$ and $F_{ss}=q^{\mathsf T}Wq$.
If $F_{ss}>0$, the information available for the state when brightness is
unknown is
\begin{equation}
 \boxed{F_{\mathrm{eff}}=F_{\theta\theta}
 -F_{\theta s}F_{ss}^{-1}F_{s\theta}.} \label{eq:schur}
\end{equation}
This expression is the \emph{Schur complement} of the brightness block.
The subtraction accounts for state changes that can be partially imitated
by a change in brightness. The following two derivations explain both its
statistical meaning and its algebraic origin.

\subsection{Derivation from local distinguishability}
Let $\eta_0=(0,s_0)$ denote the true parameters in the local state chart,
and let $P_\eta$ denote the joint distribution of all the pixel counts.
Consider the expected increase in negative log-likelihood when a candidate
parameter $\eta_0+\delta\eta$ replaces the true one:
\begin{equation}
 \mathcal D(\delta\eta)
 =\E_{\eta_0}\!\left[
 \log\frac{P_{\eta_0}(y)}{P_{\eta_0+\delta\eta}(y)}\right].
\end{equation}
This is a Kullback--Leibler divergence. Under the regularity conditions
used to define the Fisher information, its first derivative at zero
vanishes and its Hessian is $F_{\mathrm{joint}}$. Therefore
\begin{equation}
 \mathcal D(\delta\eta)
 =\frac12\delta\eta^{\mathsf T}F_{\mathrm{joint}}\delta\eta
 +o(\norm{\delta\eta}^2).
\end{equation}
The quadratic term measures how distinguishable nearby parameter values
are in the measurement statistics. With
$\delta\eta=(\delta\theta,\delta s)$, it is
\begin{equation}
 \mathcal D_{\mathrm{quad}}(\delta\theta,\delta s)
 =\frac12\left[
 \delta\theta^{\mathsf T}F_{\theta\theta}\delta\theta
 +2\delta s\,F_{s\theta}\delta\theta
 +F_{ss}(\delta s)^2\right].
 \label{eq:joint-quadratic}
\end{equation}
If brightness is known, the candidate brightness cannot change:
$\delta s=0$. The resulting state curvature is simply
$F_{\theta\theta}$.

If brightness is unknown, a candidate state change can be accompanied by
a brightness change. For a fixed $\delta\theta$, minimize
\eqref{eq:joint-quadratic} over $\delta s$ to find the candidate brightness
that makes the state change hardest to distinguish. Since $F_{ss}>0$,
this is a strictly convex scalar quadratic. Its stationary point obeys
\begin{equation}
 \frac{\partial\mathcal D_{\mathrm{quad}}}{\partial\delta s}
 =F_{s\theta}\delta\theta+F_{ss}\delta s=0,
\end{equation}
and hence
\begin{equation}
 \delta s_*=-F_{ss}^{-1}F_{s\theta}\delta\theta.
\end{equation}
Substitution gives
\begin{align}
 \min_{\delta s}\mathcal D_{\mathrm{quad}}
 &=\frac12\left[
 \delta\theta^{\mathsf T}F_{\theta\theta}\delta\theta
 -2\delta\theta^{\mathsf T}F_{\theta s}F_{ss}^{-1}
 F_{s\theta}\delta\theta
 +\delta\theta^{\mathsf T}F_{\theta s}F_{ss}^{-1}
 F_{s\theta}\delta\theta\right]\\
 &=\frac12\delta\theta^{\mathsf T}
 \left(F_{\theta\theta}-F_{\theta s}F_{ss}^{-1}F_{s\theta}\right)
 \delta\theta
 =\frac12\delta\theta^{\mathsf T}F_{\mathrm{eff}}\delta\theta.
\end{align}
Equivalently, completing the square gives the useful identity
\begin{equation}
 2\mathcal D_{\mathrm{quad}}
 =\delta\theta^{\mathsf T}F_{\mathrm{eff}}\delta\theta
 +F_{ss}\left(\delta s+
 F_{ss}^{-1}F_{s\theta}\delta\theta\right)^2.
\end{equation}
Thus $F_{\mathrm{eff}}$ is the remaining local distinguishability after
brightness has been allowed to adjust. This is a local quadratic
elimination of brightness, not a claim that the complete nonlinear
likelihood has a closed-form brightness optimizer with background present.

\subsection{Derivation from the joint covariance bound}
Suppose the state and brightness are estimated jointly. For a regular,
locally unbiased estimator, their joint covariance satisfies
\begin{equation}
 \operatorname{Cov}(\widehat\theta,\widehat s)
 \succeq F_{\mathrm{joint}}^{-1}.
\end{equation}
The relevant bound on state uncertainty is the upper-left block of this
inverse. It is not generally the inverse of the upper-left block of
$F_{\mathrm{joint}}$.

To obtain that block directly, write
$C=F_{\mathrm{joint}}^{-1}$ and use the block equations in
$F_{\mathrm{joint}}C=I$:
\begin{align}
 F_{\theta\theta}C_{\theta\theta}
 +F_{\theta s}C_{s\theta}&=I,\\
 F_{s\theta}C_{\theta\theta}+F_{ss}C_{s\theta}&=0.
\end{align}
The second equation gives
$C_{s\theta}=-F_{ss}^{-1}F_{s\theta}C_{\theta\theta}$.
Inserting this into the first equation yields
\begin{equation}
 \left(F_{\theta\theta}-F_{\theta s}F_{ss}^{-1}F_{s\theta}\right)
 C_{\theta\theta}=I.
\end{equation}
Consequently, when the necessary inverses exist,
\begin{equation}
 (F_{\mathrm{joint}}^{-1})_{\theta\theta}=F_{\mathrm{eff}}^{-1},
 \qquad
 \operatorname{Cov}(\widehat\theta)\succeq F_{\mathrm{eff}}^{-1}.
\end{equation}
Dropping the brightness row and column before inversion would give
$F_{\theta\theta}^{-1}$, the bound for known brightness, and would miss
the uncertainty due to estimating $s$.

\subsection{How much information is lost?}
The removed term is positive semidefinite, because for every real vector $a$
\begin{equation}
 a^{\mathsf T}F_{\theta s}F_{ss}^{-1}F_{s\theta}a
 =\frac{(F_{s\theta}a)^2}{F_{ss}}\ge0.
\end{equation}
Thus $F_{\mathrm{eff}}\preceq F_{\theta\theta}$: allowing unknown
brightness cannot improve the state information at fixed true parameters.
Since brightness is a scalar, the correction has rank at most one.
If $F_{\theta s}=0$, brightness and state are locally Fisher-orthogonal,
so $F_{\mathrm{eff}}=F_{\theta\theta}$ and there is no information loss.

For a single state coordinate $t$, write
\begin{equation}
 F_{\mathrm{joint}}=\begin{pmatrix}a&c\\c&b\end{pmatrix},
 \qquad F_{\mathrm{eff}}=a-\frac{c^2}{b}.
\end{equation}
Here $b$ denotes the scalar brightness-information block, not a pixel
background. If $a,b>0$, set $\gamma=c/\sqrt{ab}$; positive
semidefiniteness implies $|\gamma|\le1$. Then
$F_{\mathrm{eff}}=a(1-\gamma^2)$.
For example, $a=10$, $b=4$, and $c=6$ give $F_{\mathrm{eff}}=1$:
the state variance bound rises from $1/10$ with known brightness to $1$
with unknown brightness. At $|\gamma|=1$, the two changes are
indistinguishable to first order and the effective information vanishes.

\subsection{Projection form for computation}

A Gram matrix records inner products between a collection of vectors.
For a real matrix $G$ with columns $g_k$, its Gram matrix is
\begin{equation}
 (G^{\mathsf T}G)_{jk}=g_j^{\mathsf T}g_k.
\end{equation}
It is positive semidefinite, since
$a^{\mathsf T}G^{\mathsf T}Ga=\norm{Ga}^2\ge0$ for every real $a$.
The Fisher blocks have precisely this structure when the measurement
derivatives are weighted by their Poisson standard deviations. Define
\begin{equation}
 G=\diag(\lambda^{-1/2})sJ,\qquad
 h=\diag(\lambda^{-1/2})q.
\end{equation}
Each column of $G$ describes a weighted change in the expected measurements
along a state direction, while $h$ describes the weighted brightness change.
Consequently,
\begin{equation}
 F_{\theta\theta}=G^{\mathsf T}G,\qquad
 F_{\theta s}=G^{\mathsf T}h,\qquad F_{ss}=h^{\mathsf T}h,
\end{equation}
and the Schur complement becomes
\begin{equation}
 F_{\mathrm{eff}}=G^{\mathsf T}G
 -\frac{G^{\mathsf T}hh^{\mathsf T}G}{h^{\mathsf T}h}.
 \label{eq:gram-subtraction}
\end{equation}
Both terms in this subtraction are Gram matrices. For the second, set
$r_h=h^{\mathsf T}G/\norm{h}$; then the correction is
$r_h^{\mathsf T}r_h$. The first term contains all the local state
information for known brightness, and the second removes the contribution
that brightness changes can imitate.

\paragraph{Why subtraction can be inaccurate.}
Equation~\eqref{eq:gram-subtraction} is mathematically correct. The numerical
problem arises when the removed contribution is almost as large as the
original one: subtracting nearly equal floating-point quantities can erase
the small residual information. For a single state direction, consider
\begin{equation}
 G=\begin{pmatrix}1\\\varepsilon\end{pmatrix},\qquad
 h=\begin{pmatrix}1\\0\end{pmatrix}.
\end{equation}
In exact arithmetic,
\begin{equation}
 F_{\mathrm{eff}}=(1+\varepsilon^2)-1=\varepsilon^2.
\end{equation}
If $\varepsilon=10^{-8}$, ordinary double precision can round
$1+\varepsilon^2$ to $1$, so the computed subtraction gives zero instead
of $10^{-16}$. In a larger matrix, rounding can likewise yield spurious
small negative eigenvalues, even though the exact effective information
is positive semidefinite. These errors are particularly consequential
when the smallest eigenvalues determine rank or the inverse trace.

\paragraph{Project first, then form inner products.}
Let $P_h=hh^{\mathsf T}/(h^{\mathsf T}h)$ be the orthogonal projector
onto the brightness-derivative direction, and define
\begin{equation}
 G_\perp=(I-P_h)G
 =G-h\frac{h^{\mathsf T}G}{h^{\mathsf T}h}.
\end{equation}
Since $P_h^{\mathsf T}=P_h$ and $P_h^2=P_h$,
\begin{align}
 G_\perp^{\mathsf T}G_\perp
 &=G^{\mathsf T}(I-P_h)^{\mathsf T}(I-P_h)G\\
 &=G^{\mathsf T}(I-P_h)G
 =F_{\mathrm{eff}}.
\end{align}
Thus the effective Fisher matrix is itself a Gram matrix of the state
derivatives after removing the brightness component. In the example above,
$G_\perp=(0,\varepsilon)^{\mathsf T}$, whose squared norm directly gives
$\varepsilon^2$ without subtracting nearly equal squared quantities.

The projection still involves subtraction and cannot recover information
below the precision of the supplied derivatives, nor cure physical
near-confounding. Its advantage is that it removes the nuisance component
before squaring the derivatives and retains the Gram structure. For rank,
condition number, and inverse-trace diagnostics, prefer an SVD of
$G_\perp$ directly: forming a Gram matrix squares the condition number and
can obscure the weakest directions. The singular-value formulas in
Section~\ref{sec:diagnostics} apply with $G$ replaced by $G_\perp$ and
$F_{\mathrm{pure}}$ by $F_{\mathrm{eff}}$. A QR or SVD treatment of the
nuisance derivative span extends this
projection to several nuisance parameters.

Unknown background or calibration parameters can be handled by adding
their derivative columns and projecting out their span, provided the
nuisance model is identifiable or otherwise appropriately constrained.

\subsection{Conditioning on the number of detections}
For $b=0$, define $S=\sum_iq_i$, $p_i=q_i/S$, and
$r=J^{\mathsf T}\mathbf 1$. The probability Jacobian is
\begin{equation}
 J_p=\frac{J-pr^{\mathsf T}}{S},\qquad
 F_1=J_p^{\mathsf T}\diag(1/p_i)J_p. \label{eq:conditional}
\end{equation}
Here $F_1$ is information per detected photon. Conditional on a fixed total
of $N$ detections, the information is $NF_1$. Direct substitution gives
\begin{equation}
 F_{\mathrm{eff}}=s\left[J^{\mathsf T}\diag(1/q_i)J
 -\frac{rr^{\mathsf T}}{S}\right]=sS F_1.
\end{equation}
Thus eliminating unknown brightness in the background-free Poisson model
agrees with the conditional formulation at the expected count $sS$.

If $A^\dagger A=I$, then $S=1$ for every normalized state and $r=0$.
Brightness and state are then Fisher-orthogonal, and
$F_{\mathrm{eff}}=sF_1=sJ^{\mathsf T}\diag(1/q_i)J$.
Normalizing each column of $A$ separately does \emph{not} guarantee
$A^\dagger A=I$. Finite camera coverage can make $S$ state dependent.
With nonzero additive background, conditioning on total counts generally
does not eliminate $s$, because signal-to-background proportions depend on it.


\begin{thebibliography}{9}
\bibitem{oliveira2025}
M. Gil de Oliveira, A. L. S. Santos Junior, P. M. R. Lima,
A. C. Barbosa, B. Pinheiro da Silva, S. P\'adua, and A. Z. Khoury,
``Informationally Complete Orbital Angular Momentum Tomography with
Intensity Measurements,'' arXiv:2404.05616v3 (2025).
\url{https://arxiv.org/abs/2404.05616}.
Background for the camera POVM, Bloch-coordinate measurement matrix,
and conditional detection under obstructions.

\bibitem{li2022}
Z. Li, K. Lange, and J. A. Fessler,
``Poisson Phase Retrieval in Very Low-count Regimes,''
\emph{IEEE Transactions on Computational Imaging} (2022),
\href{https://doi.org/10.1109/TCI.2022.3209936}{doi:10.1109/TCI.2022.3209936};
arXiv:2104.00861v4.
\url{https://arxiv.org/abs/2104.00861}.
See Eqs.~(6)--(14) for the likelihood and WF step-size construction,
and Sec.~III A for spectral initialization.

\bibitem{li2016}
N. Li, C. Ferrie, J. A. Gross, A. Kalev, and C. M. Caves,
``Fisher-symmetric informationally complete measurements for pure states,''
\emph{Physical Review Letters} \textbf{116}, 180402 (2016).
\href{https://doi.org/10.1103/PhysRevLett.116.180402}{doi:10.1103/PhysRevLett.116.180402};
\url{https://arxiv.org/abs/1507.06904}.
Reference for local pure-state parameters and Fisher information,
including the distinction between local and global completeness.

\bibitem{heinosaari2013}
T. Heinosaari, L. Mazzarella, and M. M. Wolf,
``Quantum Tomography under Prior Information,''
\emph{Communications in Mathematical Physics} \textbf{318}, 355--374 (2013).
\href{https://doi.org/10.1007/s00220-013-1671-8}{doi:10.1007/s00220-013-1671-8};
\url{https://arxiv.org/abs/1109.5478}.
Reference for informational completeness restricted to pure or bounded-rank
states; the linked arXiv version includes a later correction to an outcome bound.

\bibitem{conca2015}
A. Conca, D. Edidin, M. Hering, and C. Vinzant,
``An algebraic characterization of injectivity in phase retrieval,''
\emph{Applied and Computational Harmonic Analysis} \textbf{38}, 346--356 (2015).
\href{https://doi.org/10.1016/j.acha.2014.06.005}{doi:10.1016/j.acha.2014.06.005};
\url{https://arxiv.org/abs/1312.0158}.
Reference for global complex phase-retrieval injectivity and the
rank-at-most-two Hermitian kernel criterion.

\bibitem{rasmussen2006}
C. E. Rasmussen and C. K. I. Williams,
\emph{Gaussian Processes for Machine Learning}, MIT Press (2006), Chapter 4.
\url{https://gaussianprocess.org/gpml/chapters/RW4.pdf}.
See the covariance--spectrum relation and the squared-exponential covariance
function for the interpretation of a Gaussian spatial length scale.
\end{thebibliography}
\end{document}
